\section{Exact algorithm for FVS-4CL using MSOL}
\label{subsection:exact-alg}

In this section we are looking for an algorithm that finds exact solutions to the 
FVS-4CL problem in each subgraph. For each subgraph we assume that the treewidth is bounded. 

We take a logical approach that was introduced by Courcelle~\cite{courcelle1990monadic}. 
He shows that graph properties, that can be stated in monadic second-order logic (MSOL), can be solved in 
linear time for graphs of bounded treewidth. In the following two subsections we briefly explain MSOL and 
then apply it to our problem. We follow a survey on the matter in Satzinger \cite{satzinger2014tree}. 

\subsection{MSOL - a language for graph problems}

 We want to introduce a language that translates graph problems to formulas, therefore reducing 
the problem to a satisfiability problem. Consequently, our goal is to find a formula $\Phi$ with the 
property: 

\[\text{The problem $\Pi$ is solvable for a graph $G$ \textit{iff} $G \models \Phi$, i. e. the formula 
$\Phi$ is satisfied by the graph $G$.}\]

Problem $\Pi$ is solvable in polynomial (in this case linear) time for graphs of treewidth at 
most $k$.

Monadic second-order logic is an extension of first-order logic. Additionally to having 
object variables $x, y \dots$, logical connectives $\lnot, \land, \lor, \lnot,
\rightarrow, \longleftrightarrow$ and quantifiers for object variables $\forall x, 
\exists x$, MSOL considers sets of object variables $X, Y \dots$, and relation to them 
via $\in$.

Whenever graph problems are involved, typically we express solution through 
vertices and edges. Additionally, as we mentioned, we can construct sets of vertices 
and edges. Therefore we have the following machinery:

\begin{enumerate}[label=\textbf{•}]
    \item Object variables: \textit{vertices}: $v_1, v_2, \dots$ and \textit{edges}: $e_1, e_2, \dots$. 
    \item Set variables: sets of vertices $V_1, V_2, \dots$ and sets of edges $E_1, E_2, \dots$.
    \item Binary relation $\in : \{\text{object variable}\} \times \{\text{set variable}\} \to \{0, 1\}$. Therefore 
    $v \in V$ \textit{iff} $v$ is an element of the corresponding set $V$.
    \item The $Adj(e, v_i, v_j)$ relation. It detects whether edge $e$ is an edge from vertex $v_i$ to 
    vertex $v_j$ where $v_i \neq v_j$.
\end{enumerate}

 With this machinery, one obtains a first-order language suitable for expressing relations in a graph. 
To extend this to a traditional second-order logic language, one must add quantifications over set 
variables. 

\begin{enumerate}[label=\textbf{•}]
    \item Quantification over set variables: $\forall V_i, \forall E_i$ and $\exists V_i, \exists E_i$.
\end{enumerate}

\subsection{Exact algorithm for FVS-4CL using MSOL}

Courcelle's work was extended to extended monadic second-order logic 
\cite{arnborg1991easy,borie1991deterministic,courcelle1993monadic}. These extensions allow to optimize 
over the size and weight of a free set variable.

\begin{equation}
    \label{courcelle:unweighted}
    \begin{aligned}
        \min_{F \subseteq V} & \ |F| \\
        & \land \forall v \forall u \forall w \forall z \\
        & v \in (V \backslash F) \land u \in (V \backslash F) \land w \in (V \backslash F) \land z \in (V \backslash F) \\
        & \land v \neq w \land u \neq z \\
        & \land \neg ((v, u) \in E \land (u, w) \in E \land (w, z) \in E \land (v, z) \in E)
    \end{aligned}
\end{equation}

The minimum weight FVS-4CL can be expressed as:

\begin{equation}
    \label{courcelle:weighted}
    \begin{aligned}
        \min_{F \subseteq V} & \ w(F): \\
        & \forall v \forall u \forall w \forall z \\
        & v \in (V \backslash F) \land u \in (V \backslash F) \land w \in (V \backslash F) \land z \in (V \backslash F) \\
        & \land v \neq w \land u \neq z \\
        & \land \neg((v, u) \in E \land (u, w) \in E \land (w, z) \in E \land (v, z) \in E)
    \end{aligned}
\end{equation}


\subsection{Exact algorithm for weighted FVS-BCL using MSOL}

 The minimum weight FVS-$\rho$CL can be expressed as:

\begin{equation}
    \label{courcelle:rho}
    \begin{aligned}
        \min_{F \subseteq V} & \ w(F) : \\
        & \land \forall v_1 \forall v_2 \dots \forall v_{\rho} \\
        & v_1 \in (V \backslash F) \land v_2 \in (V \backslash F) \dots \land v_{\rho} \in (V \backslash F) \\
        & \land (v_1 \neq v_2 \land v_1 \neq v_3 \land \dots \land v_1 \neq v_{\rho}) \\ 
        & \land (v_2 \neq v_1 \land v_2 \neq v_3 \land \dots \land v_2 \neq v_{\rho}) \\
        & \dots \\
        & \land (v_{\rho} \neq v_1 \land v_{\rho} \neq v_2 \land \dots \land v_{\rho} \neq v_{\rho-1}) \\
        & \land \neg((v_1, v_2) \in E \land (v_2, v_3) \in E \land \dots \land (v_{\rho-1},v_{\rho}) \in E \land (v_{\rho}, v_1) \in E)
    \end{aligned}
\end{equation}


\refstepcounter{definition}
\begin{mydefinition}[Courcelle's theorem \cite{courcelle1990monadic}]{Theorem}
    \label{theorem:courcelle}

    Assume that $\phi$ is a MSOL formula and $G$ is an n-vertex graph, and there is an evaluation 
    of all of the free variables of $\phi$. Suppose, that a tree decomposition of $G$ of width $t$ 
    is provided. Then there exists an algorithm that verifies that $\phi$ is satisfied 
    in $G$ in time $f(||\phi||, t) \cdot n$, for some computable function $f$. 
    
\end{mydefinition}

\refstepcounter{definition}
\begin{mydefinition}[MSOL Solvability of FVS-$\rho$CL on Bounded-Treewidth Graphs]{Corollary}
    \label{cor:msol-fvs-rho}
    Let $\rho\ge 3$ be a constant, and let $G=(V,E)$ be an undirected graph of treewidth at 
    most $\tw$, equipped with a vertex-weight function $w:V\to\mathbb{N}$.  Then the 
    minimum-weight feedback-vertex-set that breaks all $\rho$-cycles in $G$ 
    (i.e.\ solves FVS-$\rho$CL) can be computed in time
    \[
        f(\rho, \tw) \cdot n^{\mathcal{O}(1)},
    \]
    where $f$ is some computable function depending only on $\rho$ and $\tw$. 
\end{mydefinition}
  

\subsection{Exact algorithm for FVS-4CL using dynamic programming}

\label{alg:dp}

In order to find an algorithm that solves FVS-4CL, one has to define the following notions. 
For a more thorough explanation of those concepts, see Bodlaender~\cite[Section 4.]{bodlaender1997treewidth}.

\begin{enumerate}
    \item Define the notion of a \textit{solution}. In our case for the FVS-4CL problem, the solution for a graph 
    $G$ is a set $F$, such that $G - F$ does not contain any cycles of length four.

    \item Define the notion of a \textit{partial solution}.
    For the subgraph $G_i = (V_i, E_i)$, partial solutions $F_i$ are restrictions of solutions 
    given by the subsets $F_i \subseteq V_i$. 
    
    \item Define the notion of \textit{extension of a partial solution}. A solution $F$ is an extension 
    of a partial solution $F_i$ for $G_i$ if $F \cap V_i = F_i$.

    \item Define the notion of a \textit{characteristic} of a partial solution. 
    For the set $X_i$ the vertices are partitioned twofold.

    \begin{enumerate}
        \item The class $I \subseteq X_i$ corresponds to all vertices that are in the partial solution 
        for the current node $X_i$.
        \item The class $\mathcal{F} \subseteq X_i \times X_i$ is comprised of all of the pairs of 
        vertices $v,u \in X_i: 
        (v, u) \in \Phi$ if they have a common neighbour $w \in V_i \backslash X_i$, such that 
        $w \notin F_i$.
    \end{enumerate}
    
    \[ch(G_i, F_i) = (\{v \in X_i: v \in F_i\}, 
    \{(v, u) \in X_i \times X_i: \exists w \in V_i \backslash X_i, 
    (v, w) \in E_i, (u, w) \in E_i, w \notin F_i\})\]

    The valuation of a characteristic $ch(G_i, F_i)$ is a table $c[i, I, \mathcal{F}]$.
    It represents the minimum weight 
    of a partial solution $F_i$ in $G_i$ and $c[i, I, \mathcal{F}] \in \mathbb{N} \cup \{\infty\}$.
    The table keeps the value of infinity if no such set exists.

    \[c[i, I, \mathcal{F}] = \min\{w(W) : W \text{ is a FVS-4CL of } G_i \land W \cap X_i = I\}\] 

    \item Define the notion of a \textit{full set of characteristics}.

    The full set for a node $X_i$ is a valuation for each characteristic $I \in \{0, 1\}^{|X_i|}$ 
    and $\mathcal{F} \in \{0, 1\}^{X_i \times X_i}$. Since there are at most $\tw+1$ elements, 
    there are at most $2^{(\tw+1)} \cdot 2^{\binom{\tw+1}{2}}$ entries. That yields a storage 
    capacity of $2^{\binom{\tw+1}{2} + (\tw+1)} = 2^{\mathcal{O}(\tw^2)}$. 

\end{enumerate}

\textbf{Leaf node:}

For a leaf node $i$, we have that $X_i = \emptyset$. Hence 

\[c[i, \emptyset, \emptyset] = 0\]

\textbf{Introduce node:}

Let $i$ be an introduce node with child node $j$. We have that $X_i = X_j \cup \{v\}$ for some 
vertex $v \notin X_j$. There are two possible values.
\[
\begin{aligned}
    c[i, I \cup \{v\}, \mathcal{F}] &= w(v) + c[j, I, \mathcal{F}] \\
    c[i, I, \mathcal{F}] &= 
    \begin{cases}
        \infty, & \text{if } \exists u, w \in X_i: (u, w) \in \mathcal{F}, u \notin I, w \notin I, (v, u) \in E_i, (v, w) \in E_i \\
        c[j, I, \mathcal{F}], & \text{otherwise}
    \end{cases}
\end{aligned}
\]

\textbf{Forget node:}

Let $i$ be a forget node with child node $j$. We have that $X_i = X_j \backslash \{v\}$ for 
some vertex $v \in X_j$.

\[c[i, I, \mathcal{F}] = \min(c[j, I \cup \{v\}, \mathcal{F}],
c[j, I, \mathcal{F} \cup \{(u, w): (v,u) \in E_i, (v, w) \in E_i\}])\]

\textbf{Join node:}

Let $i$ be a join node with children $j_1$ and $j_2$. Recall that $X_i = X_{j_1} = X_{j_2}$. 
Let $\mathcal{F}_1$ and $\mathcal{F}_2$ be the $\mathcal{F}$ set values of $X_{j_1}$ and $X_{j_2}$ respectively.

\[
\begin{aligned}
    c[i, I, \mathcal{F}_1 \cup \mathcal{F}_2] =
    \begin{cases}
        \infty, & \text{if } \exists u, w \in X_i : (u, w) \in \mathcal{F}_1, \\
                & \quad (u, w) \in \mathcal{F}_2, u \notin I, w \notin I \\
        c[j_1, I, \mathcal{F}_1] + c[j_2, I, \mathcal{F}_2] - w(I), & \text{otherwise}
    \end{cases}
\end{aligned}
\]

